\section{Human Evaluation Results}
\label{sec:results-human}

The study closed at 27 complete raters, above the 18--20 recruitment target,
giving 162 pairwise comparisons. Every rater who finished all six tasks is
included; one incomplete session is excluded, as specified before data
collection. Participants self-reported as students~(10), junior
developers~(7), mid-level developers~(8), and senior developers~(2), with
programming experience ranging from under one year~(3) through one to
three~(9), three to five~(12), and five to ten years~(3); 19 of 27
reported reviewing pull requests at least weekly.

\subsection{Primary Endpoint}
\label{sec:results-human-primary}

On H1 (which review better names the specific functions, classes, or files
the change could affect), the mean rater preference for \kg{} is $0.710$,
with a bootstrap interval of $[0.623, 0.787]$ excluding the no-preference
value of~$0.5$. Wilcoxon signed-rank $p = 0.0005$ against~$0.5$, with
rank-biserial $r = 0.80$. Across the 162 comparisons, raters picked the
\kg{} review 98 times and the \baseline{} review 30 times, with 29 ties and
5 ``neither'' judgements.

\subsection{Overall Usefulness}
\label{sec:results-human-overall}

Mean preference on overall usefulness is $0.512$, $[0.423, 0.605]$.
The interval includes $0.5$, so the data are consistent with no preference.
This question was set aside for estimation before the study started, so no
significance test is reported. Raters preferred the graph review when asked
which review names the affected code, and they showed no preference when asked
which review was more useful overall. The result is therefore specific to
affected-code identification.

\subsection{Exploratory Boundary}

Among the five exploratory criteria, only code clarity remains distinguishable
after correction, and it favours the \baseline{} review. The graph review is
ahead only on the primary affected-code question.
Figure~\ref{fig:human-study} summarises all criteria on one axis;
Table~\ref{tab:human-exploratory} (Appendix) reports the exploratory
estimates and Holm-corrected $p$-values.

\begin{figure}[htbp]
    \centering
    % !TeX root = ../thesis.tex
% GENERATED FILE - do not edit by hand; your changes will be lost.
% Produced by scripts/generate_thesis_figures.py
% Regenerate: python3 scripts/generate_thesis_figures.py
% Computed from:
%   experiments/2026-07-06_user_study_prs/RESULTS_HUMAN_V4.json (n = 27 raters, 6 pull requests each)
%
\begin{tikzpicture}
\node[anchor=west,font=\footnotesize] at (-5.6000,0.0000) {\emph{Primary (confirmatory)}};
\node[anchor=west,font=\footnotesize] at (-5.6000,-0.4800) {\textbf{H1} names affected code};
\draw[thin,black!55] (-5.6000,-0.7104) -- (6.8500,-0.7104);
\node[anchor=west,font=\footnotesize] at (-5.6000,-0.9600) {\emph{Secondary (estimation only)}};
\node[anchor=west,font=\footnotesize] at (-5.6000,-1.4400) {\textbf{Overall} usefulness};
\draw[thin,black!55] (-5.6000,-1.6704) -- (6.8500,-1.6704);
\node[anchor=west,font=\footnotesize] at (-5.6000,-1.9200) {\emph{Exploratory (Holm-corrected)}};
\node[anchor=west,font=\footnotesize] at (-5.6000,-2.4000) {\textbf{H2} describes failure situations};
\node[anchor=west,font=\footnotesize] at (-5.6000,-2.8800) {\textbf{H3} points to test files};
\node[anchor=west,font=\footnotesize] at (-5.6000,-3.3600) {\textbf{H4} explains why something is a problem};
\node[anchor=west,font=\footnotesize] at (-5.6000,-3.8400) {\textbf{H5} comments on code clarity};
\node[anchor=west,font=\footnotesize] at (-5.6000,-4.3200) {\textbf{H6} covers the important issues};
\draw[densely dashed,black!55] (1.9304,-0.2640) -- (1.9304,-4.5360);
\draw[kgblue,line width=0.9pt] (3.4159,-0.4800) -- (5.3000,-0.4800);
\draw[kgblue,line width=0.9pt] (3.4159,-0.5500) -- (3.4159,-0.4100);
\draw[kgblue,line width=0.9pt] (5.3000,-0.5500) -- (5.3000,-0.4100);
\draw[fill=kgblue,draw=kgblue] (4.3942,-0.4800) circle (2.000pt);
\node[anchor=east,font=\scriptsize] at (6.8500,-0.4800) {0.0005};
\draw[modegrey,line width=0.9pt] (1.0246,-1.4400) -- (3.1623,-1.4400);
\draw[modegrey,line width=0.9pt] (1.0246,-1.5100) -- (1.0246,-1.3700);
\draw[modegrey,line width=0.9pt] (3.1623,-1.5100) -- (3.1623,-1.3700);
\draw[fill=white,draw=modegrey,line width=0.6pt] (2.0754,-1.4400) circle (2.000pt);
\node[anchor=east,font=\scriptsize] at (6.8500,-1.4400) {---};
\draw[modegrey,line width=0.9pt] (0.9159,-2.4000) -- (2.2565,-2.4000);
\draw[modegrey,line width=0.9pt] (0.9159,-2.4700) -- (0.9159,-2.3300);
\draw[modegrey,line width=0.9pt] (2.2565,-2.4700) -- (2.2565,-2.3300);
\draw[fill=white,draw=modegrey,line width=0.6pt] (1.5681,-2.4000) circle (2.000pt);
\node[anchor=east,font=\scriptsize] at (6.8500,-2.4000) {0.499};
\draw[modegrey,line width=0.9pt] (0.6986,-2.8800) -- (2.0754,-2.8800);
\draw[modegrey,line width=0.9pt] (0.6986,-2.9500) -- (0.6986,-2.8100);
\draw[modegrey,line width=0.9pt] (2.0754,-2.9500) -- (2.0754,-2.8100);
\draw[fill=white,draw=modegrey,line width=0.6pt] (1.3870,-2.8800) circle (2.000pt);
\node[anchor=east,font=\scriptsize] at (6.8500,-2.8800) {0.499};
\draw[modegrey,line width=0.9pt] (1.2058,-3.3600) -- (2.4377,-3.3600);
\draw[modegrey,line width=0.9pt] (1.2058,-3.4300) -- (1.2058,-3.2900);
\draw[modegrey,line width=0.9pt] (2.4377,-3.4300) -- (2.4377,-3.2900);
\draw[fill=white,draw=modegrey,line width=0.6pt] (1.8217,-3.3600) circle (2.000pt);
\node[anchor=east,font=\scriptsize] at (6.8500,-3.3600) {1.000};
\draw[modegrey,line width=0.9pt] (0.3000,-3.8400) -- (1.6768,-3.8400);
\draw[modegrey,line width=0.9pt] (0.3000,-3.9100) -- (0.3000,-3.7700);
\draw[modegrey,line width=0.9pt] (1.6768,-3.9100) -- (1.6768,-3.7700);
\draw[fill=modegrey,draw=modegrey] (0.9159,-3.8400) circle (2.000pt);
\node[anchor=east,font=\scriptsize] at (6.8500,-3.8400) {0.028};
\draw[modegrey,line width=0.9pt] (0.9522,-4.3200) -- (2.5101,-4.3200);
\draw[modegrey,line width=0.9pt] (0.9522,-4.3900) -- (0.9522,-4.2500);
\draw[modegrey,line width=0.9pt] (2.5101,-4.3900) -- (2.5101,-4.2500);
\draw[fill=white,draw=modegrey,line width=0.6pt] (1.7130,-4.3200) circle (2.000pt);
\node[anchor=east,font=\scriptsize] at (6.8500,-4.3200) {1.000};
\draw[thin,black!55] (0.0000,-4.7760) -- (5.6000,-4.7760);
\draw[thin,black!55] (0.7565,-4.7760) -- (0.7565,-4.8560);
\node[anchor=center,font=\scriptsize] at (0.7565,-5.0560) {0.4};
\draw[thin,black!55] (1.9304,-4.7760) -- (1.9304,-4.8560);
\node[anchor=center,font=\scriptsize] at (1.9304,-5.0560) {0.5};
\draw[thin,black!55] (3.1043,-4.7760) -- (3.1043,-4.8560);
\node[anchor=center,font=\scriptsize] at (3.1043,-5.0560) {0.6};
\draw[thin,black!55] (4.2783,-4.7760) -- (4.2783,-4.8560);
\node[anchor=center,font=\scriptsize] at (4.2783,-5.0560) {0.7};
\draw[thin,black!55] (5.4522,-4.7760) -- (5.4522,-4.8560);
\node[anchor=center,font=\scriptsize] at (5.4522,-5.0560) {0.8};
\node[anchor=center,font=\footnotesize] at (2.8000,-5.4960) {Rater preference for the \kg{} review ($0.5 =$ no preference)};
\node[anchor=east,font=\scriptsize] at (6.8500,0.0000) {$p$};
\end{tikzpicture}

    \caption{Rater preference for the \kg{} review by criterion ($n = 27$).
             Bars are bootstrap intervals; the dashed line is no preference.
             Filled markers indicate significance under the governing test.}
    \label{fig:human-study}
\end{figure}
